\subsection{Rule C}

If we have some wf of the form $(\exists x) \B(x)$, we can choose some $b$ such that $\B(b)$. This can often simplify proofs, as this is basically a shortcut for a longer sequence of steps. Check page 78-79 for an example. There are some rules, however. A rule C deduction in a FOT K is defined in the following manner:

$\Gamma \vdash_\text{C} \B$ iff there is a sequence of wfs $\D_1, \dots, \D_n$ such that $\D_n$ is $\B$ and the following conditions hold:

\begin{enumerate}
    \item For each $i<n$, either
    \begin{enumerate} [label = \alph*.]
        \item $\D_i$ is an axiom of K
        \item $\D_i$ is in $\Gamma$
        \item $\D_i$ follows by MP or Gen from preceding wfs in the sequence
        \item There is a preceding wf $(\exists x) \C(x)$ such that $\D_i$ is $\C(d)$, where $d$ is a \textbf{new} individual constant (Rule C).
    \end{enumerate}
    \item As axioms in condition 1a, we also can use all logical axioms that involve the new individual constants already introduced in the sequence by rule C.
    \item No application of Gen is made using a variable that is free in some $(\exists x) \C(x)$ to which rule C has been previously applied.
    \item $\B$ contains none of the new individual constants introduced in the sequence in any application of rule C.
\end{enumerate}

Condition 3 is important because if an application of rule C to a wf $(\exists x)\C(x)$ 
yields $\C(d)$, then the object referred to by $d$ may depend on the values of the free variables in $(\exists x)\C(x)$. So that one object may not satisfy $\C(x)$ for \textit{all} values of the free variables in $(\exists x)\C(x)$. There is a good example on page 80 which compares the wfs $(\forall x)(\exists y) A_1^2(x,y)$ and $(\exists y)(\forall x) A_1^2(x,y)$, where there is some interpretation which makes the former true but the latter false.

\textbf{Proposition 2.10}
If $\Gamma \vdash_\text{C} \B$, then $\Gamma \vdash \B$. Moreover, from the following proof one can verify that, if there is an application of Gen in the new  proof of  $\B$ using $\Gamma$ using a certain variable and applied to a wf depending upon a certain wf of $\Gamma$, then there was such an application of Gen in the original proof.

\begin{enumerate}
    \item Let $(\exists y_1)\C_1(y_1), \dots, (\exists y_k)\C_k(y_k)$ be the wfs in order of occurrence to which rule C is applied in the proof of $\Gamma \vdash_\text{c} \B$, and let $d_1, \dots, d_k$ be the corresponding new individual constants. 
    \item Then, $\Gamma, \C_1(d_1), \dots, \C_k(d_k) \vdash \B$. 
    \item Now, by condition 3 of the definition above, Corollary 2.6 is now applicable, yielding $\Gamma, \C_1(d_1), \dots, \C_{k-1}(d_{k-1}) \vdash \C_k(d_k) \imp \B$.
    \item We replace $d_k$ everywhere by a variable $z$ that does not occur in the proof.
    \item So:
    \[\Gamma, \C_1(d_1), \dots, \C_{k-1}(d_{k-1}) \vdash \C_k(z) \imp \B\]
    \item Now, by Gen:
    \[\Gamma, \C_1(d_1), \dots, \C_{k-1}(d_{k-1}) \vdash (\forall x)(\C_k(z) \imp \B)\]
    \item So, by Exercise 2.32(d):
    \[\Gamma, \C_1(d_1), \dots, \C_{k-1}(d_{k-1}) \vdash (\exists y_k)\C_k(y_k) \imp \B\]
    \item Notice that we can obtain:
    \[\Gamma, \C_1(d_1), \dots, \C_{k-1}(d_{k-1}) \vdash (\exists y_k)\C_k(y_k)\]
    \item So, by MP we obtain:
    \[\Gamma, \C_1(d_1), \dots, \C_{k-1}(d_{k-1}) \vdash \B\]
\end{enumerate}
This argument can be repeated to eliminate $\C_{k-1}(d_{k-1}), \dots, \C_1(d_1)$ to obtain $\Gamma \vdash \B$. I don't know why Mendelson didn't just use regular induction for this, this is a strange proof relative to the others in this chapter.
